% =====================================================================
\section{How the rhythm breaks}\label{sec:bifur}
A separate question, posed independently of the redox part: what happens to the alternation of sleep and waking as damage inhibits the sleep nucleus ever more strongly. The expectation was this: on approach to a critical value chaos should appear --- that is, the circadian rhythm breaks \emph{while} the 24-hour external pacemaker is preserved. This is cheap to compute: the switch is deterministic, and neither the cells nor the redox part are needed at all.

The right observable here is the \textbf{rotation number}: how many days of the external forcing fall on one cycle of sleep and waking. The rhythm does not break smoothly: the rotation number stands at rational values on finite intervals of the parameter and jumps between them (fig.~\ref{fig:slow}c):
\[1\!:\!1\;\to\;5\!:\!4\;\to\;3\!:\!2\;\to\;2\!:\!1\;\to\;3\!:\!1\;\to\;4\!:\!1\;\to\;\dots\;\to\;9\!:\!1\;\to\;\text{no sleep}.\]
The transitions fall at $kD\approx4.77$; 4.78; 4.79; 5.09; 5.16 and then ever more densely: five steps (5:1, 6:1, 7:1, 8:1, 9:1) fit between 5.172 and 5.178\,mV, after which sleep disappears altogether.

A separate run with a fine scan and refinement of the boundaries (fig.~\ref{fig:stairs}) shows that there are far more steps: \textbf{52 over an interval of 0.9\,mV, and in 34 of them the cycle lengths alternate}. Between the integer lockings stands a dense accumulation of fractional ones --- 5:4, 6:5, 4:3, 13:10, 11:8, 9:7, 7:5, 10:7, 16:11, 8:5, 12:7 --- and between those, points where the pattern does not repeat at all. The step widths fall with the denominator by many orders of magnitude: for 2:1 it is 0.3\,mV, for 3:2 about $8\cdot10^{-3}$, for 11:8 already $5\cdot10^{-6}$. A uniform scan cannot catch this: only the boundaries can be refined.

\figp{fig_stairs.png}{\label{fig:stairs}The devil's staircase in full, from a separate run with refinement of the boundaries (about 60 cycles per point after the transient was discarded). \textbf{Above} --- the whole interval: the flat steps of the integer lockings 1:1, 2:1, 3:1 and on to 9:1 packed together just before sleep disappears. \textbf{Below} --- a close-up of the region with the most steps: in orange, lockings with \emph{alternating} cycle lengths; in grey, regimes with no repeating pattern. The inset shows the pattern of the 3:2 locking itself: a long cycle of 47\,h, a short one of 25\,h, and so round again.}

\textbf{This is a locking staircase, not a period-doubling cascade.} In a doubling cascade the period is multiplied by two and takes no other values, and the Feigenbaum constant describes how the steps converge. Here the period adds one missed night at a time, and fractional steps stand between the integer ones --- these are Arnold tongues, the ordinary picture of an oscillation entrained by an external force. To invoke the Feigenbaum constant here would be an error.

\textbf{The most interesting thing is the fractional steps.} In them the cycle lengths are not equal: they follow a strict pattern and repeat. At $kD=4.772$\,mV the pattern has four cycles over five days: 25, 24, 24, 47\,h --- three nearly ordinary days, then one doubled, and again (the inset in fig.~\ref{fig:slow}c). At 4.779 the pattern has two: 47 and 25\,h. At 4.783, three: 46, 26, 24\,h. Checked for stability: the pattern is the same at 120 and at 400 days of computation, so it is not a transient.

This matters for the following reason. Such behaviour is not ``no rhythm'' but \emph{another rhythm}, with a period of its own of several days, and this is exactly how part of the clinical picture of disturbed sleep looks. The model produces it by itself, without any noise and without a single added equation --- simply as a consequence of entrainment. A prediction written down before the computation: with noise added, these windows should not disappear but blur into an \emph{intermittent} alternation --- the pattern holds, then breaks down, then holds again.

\textbf{The principal limitation is honest and unpleasant --- but it is not the one that stood here in version~4.} The ceiling is set not by the damage but by the coupling itself: identically $kD\le k$, so at $k=2$\,mV the step at 4.77\,mV is unattainable \emph{for any} damage. A coupling of no less than 4.77\,mV is needed, that is, 2.4 times stronger than the working one. At $k=5$\,mV the model \emph{does} reach the staircase: in sets that have fallen into the irreversible state, $kD=4.93$ was measured at a damaged fraction of 0.999 and waking of 0.974. So the limitation is not in the mechanism but in the calibration of a single constant.

\textbf{The previous edition named two ways out; the second is struck off.} ``Chronic damage holding the fraction of broken cells high'' does not work, and this is evident without any computation: in the broken sets the fraction is already at its ceiling, and the ceiling equals $k$. One way out remains --- the strength of the ``damage--sleep'' coupling. The earlier numbers in this paragraph (``twofold'', ``fourfold'') were wrong and did not even agree with each other; their root is analysed in section~\ref{sec:prec} as the seventh statistical trap.

\begin{note}
The check written down before the computation asked the wrong thing: it looked among the periods for the values 1, 2 and 4 --- and found them, but would have found them in a locking staircase as well. The verdict is left as it stands, and beside it are recorded the two parts that were missing: is there locking at a period that is not a power of two, and are there regimes with a pattern longer than one cycle? Separately, the scan grid had to be corrected: the 5:4 and 3:2 steps have a width of about 0.005\,mV and are invisible on a coarse grid --- the first run saw 1, 2, 4 and not a single fractional one.
\end{note}

\begin{ineq}
Why a staircase and not a cascade was what had to come out here. The switch is a relaxation oscillator with an intrinsic period of about 25\,h; the daily term is an external force of period 24\,h. For such a pair the ratio of the periods is the rotation number, and it locks at rational $n\!:\!m$ on intervals of finite width (Arnold tongues), varying between them along a devil's staircase. The width of a tongue falls rapidly with the denominator $m$: for 2:1 it is about 0.3\,mV, for 3:2 and 5:4 about 0.005, that is, sixty times narrower. A staircase can be distinguished from a doubling cascade by exactly one sign: the presence of locking at a period that is not a power of two.
\end{ineq}
